\chapter{Problem i wykorzystane technologie}
\label{ch:problem}

\section{Równanie ciepła i metoda różnic skończonych}
\label{sec:heat-fdm}

\subsection{Model fizyczny}

Rozważanym zjawiskiem fizycznym jest dyfuzja ciepła w~jednorodnym, izotropowym
ośrodku wypełniającym sześcian jednostkowy \(\Omega = [0,1]^3\). Pole skalarne
\(u(\mathbf{x},t)\) opisuje temperaturę w~punkcie \(\mathbf{x}=(x,y,z)\in\Omega\)
w~chwili \(t\ge 0\). Ewolucję tego pola opisuje równanie różniczkowe
cząstkowe znane jako równanie przewodnictwa ciepła, zwane też równaniem dyfuzji:
\begin{equation}
  \frac{\partial u}{\partial t} \;=\; \alpha\,\nabla^2 u
  \;=\; \alpha\!\left(
    \frac{\partial^2 u}{\partial x^2}
    + \frac{\partial^2 u}{\partial y^2}
    + \frac{\partial^2 u}{\partial z^2}
  \right),
  \label{eq:heat3d}
\end{equation}
gdzie \(\nabla^2\) oznacza operator Laplace'a, natomiast
\(\alpha>0\) to współczynnik dyfuzyjności cieplnej (ang.\ \emph{thermal
diffusivity}), wyrażony w~jednostkach \si{\metre\squared\per\second}. Wielkość
\(\alpha\) łączy przewodność cieplną ośrodka z~jego pojemnością cieplną i~gęstością;
im większa jej wartość, tym szybciej lokalne różnice temperatury ulegają
wygładzeniu. W~zaimplementowanym programie przyjęto bezwymiarową wartość
\(\alpha = \num{0.25}\) jako parametr konfiguracyjny.

\subsection{Warunki początkowe i brzegowe}

Rozwiązanie równania~\eqref{eq:heat3d} jest jednoznaczne dopiero po~zadaniu
warunku początkowego oraz warunków brzegowych. W~zaimplementowanym rozwiązaniu
warunek początkowy przyjmuje jednorodne pole \(u(\mathbf{x},0)=1\) w~całej
dziedzinie, w~którym umieszczono nagrzaną warstwę 
o~grubości \texttt{hotThickness} płaszczyzn siatki przylegających do~ściany
\(y=0\), o~podwyższonej wartości \(u=2\). Warstwa ta modeluje gorącą ścianę,
od~której ciepło rozchodzi się w~głąb sześcianu.

Na~brzegu dziedziny przyjęto warunki Dirichleta o~zamrożonych wartościach: cała
zewnętrzna warstwa komórek siatki nigdy nie jest aktualizowana i~zachowuje swoje
wartości początkowe przez cały czas symulacji. W~praktyce oznacza to,
że aktualizowane jest wyłącznie wnętrze dziedziny, a~warstwa gorąca przy
\(y=0\) pozostaje utrzymywana na~poziomie \(u=2\), pełniąc rolę stałego źródła
ciepła. Pozostałe ściany, utrzymywane na~poziomie
\(u=1\), działają jak zbiorniki ciepła o~ustalonej temperaturze. Taki dobór
warunków prowadzi do~nietrywialnego, monotonicznie ewoluującego pola, które
w~granicy czasu dąży do~stanu ustalonego z~gradientem rozpiętym pomiędzy gorącą
a~zimnymi ścianami. \Cref{fig:diffusion} ilustruje jakościowo charakter tej
ewolucji na~dwuwymiarowym przykładzie: podtrzymywane źródło ciepła stopniowo
wygładza swoje ostre szczegóły, a~ciepło rozchodzi się w~otoczenie.

\begin{figure}[htbp]
  \centering
  \includegraphics[width=\linewidth]{diffusion_2d.png}
  \caption[Ewolucja pola dyfuzji 2D]{Ewolucja pola temperatury w~dwuwymiarowej
  symulacji dyfuzji, wykonana dwuwymiarowym programem Chapel (dziedzina
  \num{160}$\times$\num{160}). Od~lewej: początkowe, podtrzymywane źródło ciepła
  w~kształcie grzebienia, a~następnie stopniowe wygładzanie ostrych szczegółów
  i~propagacja ciepła w~otoczenie, aż do~ukształtowania się gładkiego gradientu.
  Barwa odwzorowuje temperaturę: błękit oznacza zimno, czerwień gorąco. Ta
  sama zasada obowiązuje w~przypadku trójwymiarowym rozważanym w~pracy.}
  \label{fig:diffusion}
\end{figure}

\subsection{Dyskretyzacja: jawna metoda różnic skończonych}

Do~rozwiązania równania~\eqref{eq:heat3d} zastosowano jawną metodę różnic
skończonych (ang.\ \emph{Finite Difference Method}, FDM) z~jawnym schematem
Eulera w~czasie i~centralnymi różnicami drugiego rzędu w~przestrzeni~\cite{leveque,strikwerda}. Dziedzinę
ciągłą zastępuje się regularną siatką kartezjańską o~\(n_x\times n_y\times n_z\)
węzłach. Krok siatki w~każdym z~kierunków dobrany jest tak, aby węzły skrajne
pokrywały się z~brzegiem sześcianu jednostkowego:
\begin{equation}
  \Delta x = \frac{1}{n_x-1}, \qquad
  \Delta y = \frac{1}{n_y-1}, \qquad
  \Delta z = \frac{1}{n_z-1}.
  \label{eq:spacing}
\end{equation}
Wartość pola w~węźle \((i,j,k)\) w~kroku czasowym \(n\) oznacza się przez
\(u_{i,j,k}^{\,n}\). Do~przybliżenia pierwszej pochodnej czasowej użyto ilorazu
różnicowego przedniego, a~drugich pochodnych przestrzennych standardowych
trzypunktowych ilorazów centralnych, na~przykład dla~kierunku
\(x\):
\begin{equation}
  \frac{\partial^2 u}{\partial x^2}(x_i,y_j,z_k)
  \approx
  \frac{u_{i-1,j,k}^{\,n} - 2\,u_{i,j,k}^{\,n} + u_{i+1,j,k}^{\,n}}{\Delta x^2},
  \label{eq:central}
\end{equation}
i~analogicznie dla~kierunków \(y\) oraz \(z\). Podstawienie tych przybliżeń
do~\eqref{eq:heat3d} i~rozwiązanie względem wartości w~następnym kroku prowadzi
do~jawnej formuły aktualizacji.

Dla~zwięzłości wprowadza się współczynniki wagowe zależne od~osi, łączące
dyfuzyjność, krok czasowy oraz krok siatki:
\begin{equation}
  a_x = \frac{\alpha\,\Delta t}{\Delta x^2}, \qquad
  a_y = \frac{\alpha\,\Delta t}{\Delta y^2}, \qquad
  a_z = \frac{\alpha\,\Delta t}{\Delta z^2}.
  \label{eq:weights}
\end{equation}
Przy tych oznaczeniach reguła aktualizacji przyjmuje zwartą postać
siedmiopunktowego stencila (ang.\ \emph{7-point stencil}):
\begin{align}
  u_{i,j,k}^{\,n+1} =\; & u_{i,j,k}^{\,n}
    + a_x\!\left(u_{i-1,j,k}^{\,n} - 2\,u_{i,j,k}^{\,n} + u_{i+1,j,k}^{\,n}\right) \notag\\
    & + a_y\!\left(u_{i,j-1,k}^{\,n} - 2\,u_{i,j,k}^{\,n} + u_{i,j+1,k}^{\,n}\right) \notag\\
    & + a_z\!\left(u_{i,j,k-1}^{\,n} - 2\,u_{i,j,k}^{\,n} + u_{i,j,k+1}^{\,n}\right).
  \label{eq:update}
\end{align}
Nowa wartość w~każdym węźle zależy wyłącznie od~wartości w~tym węźle oraz
w~jego sześciu bezpośrednich sąsiadach z~poprzedniego kroku czasowego: dwóch
wzdłuż każdej z~trzech osi. Ten wzorzec dostępu do~pamięci, komórka centralna
plus sześć sąsiadów, przedstawiono na~\cref{fig:stencil}. Zależność od~stanu
wyłącznie z~poprzedniego kroku jest cechą schematu jawnego i~czyni każdy krok
czasowy w~pełni równoległym względem węzłów siatki.

\begin{figure}[htbp]
  \centering
  \begin{tikzpicture}[scale=1.0,
      cell/.style={circle,draw,thick,minimum size=7mm,inner sep=0pt},
      center/.style={circle,draw,thick,fill=codekw!18,minimum size=7.5mm,inner sep=0pt},
      nb/.style={circle,draw,thick,fill=codestr!12,minimum size=7mm,inner sep=0pt},
      arr/.style={-{Latex[length=2mm]},thick,gray}]
    % central node
    \node[center] (C) at (0,0) {\small $u_{i,j,k}$};
    % x-neighbours
    \node[nb] (Xm) at (-2.6,0) {\small $u_{i\text{-}1}$};
    \node[nb] (Xp) at ( 2.6,0) {\small $u_{i\text{+}1}$};
    % y-neighbours (vertical)
    \node[nb] (Ym) at (0,-2.2) {\small $u_{j\text{-}1}$};
    \node[nb] (Yp) at (0, 2.2) {\small $u_{j\text{+}1}$};
    % z-neighbours (diagonal, pseudo-3D)
    \node[nb] (Zm) at (-1.5,-1.25) {\small $u_{k\text{-}1}$};
    \node[nb] (Zp) at ( 1.5, 1.25) {\small $u_{k\text{+}1}$};
    % connections
    \foreach \n in {Xm,Xp,Ym,Yp,Zm,Zp} \draw[arr] (\n) -- (C);
    % axis labels
    \node[gray] at (2.6,-0.55) {\footnotesize oś $x$};
    \node[gray] at (0.75,2.2) {\footnotesize oś $y$};
    \node[gray] at (1.9,0.55) {\footnotesize oś $z$};
  \end{tikzpicture}
  \caption[Stencil siedmiopunktowy 3D]{Siedmiopunktowy stencil trójwymiarowy.
  Nowa wartość w~komórce centralnej \(u_{i,j,k}\), zaznaczonej na~niebiesko, jest ważoną
  kombinacją jej wartości bieżącej oraz sześciu bezpośrednich sąsiadów
  zaznaczonych na~czerwono, wzdłuż osi \(x\), \(y\) i~\(z\), zgodnie z~równaniem~\eqref{eq:update}.}
  \label{fig:stencil}
\end{figure}

\subsection{Stabilność numeryczna i warunek CFL}
\label{sec:cfl}

Zaletą schematów jawnych jest prostota i~niski koszt pojedynczego kroku: nie wymagają one rozwiązywania układów równań, w~przeciwieństwie do~schematów
niejawnych. Ceną za~tę prostotę jest jednak warunkowa stabilność: krok czasowy
\(\Delta t\) nie może być dowolnie duży~\cite{leveque,morton-mayers}. Zbyt duży krok powoduje, że błędy
zaokrągleń i~składowe wysokoczęstotliwościowe rozwiązania numerycznego są
wzmacniane w~każdej iteracji zamiast tłumione, co prowadzi do~gwałtownej,
niefizycznej rozbieżności rozwiązania.

Analiza stabilności metodą von~Neumanna~\cite{strikwerda,leveque} dla~schematu
jawnego zastosowanego do~równania dyfuzji w~trzech wymiarach prowadzi do~warunku
Couranta--Friedrichsa--Lewy'ego (CFL). Dla~liniowego schematu o~stałych
współczynnikach współczynnik wzmocnienia pojedynczego modu Fouriera wynosi
\(g = 1 - 4a_x\sin^2(\theta_x/2) - 4a_y\sin^2(\theta_y/2) - 4a_z\sin^2(\theta_z/2)\),
a~żądanie \(\lvert g\rvert\le 1\) dla~najbardziej niekorzystnego modu o~naprzemiennym
znaku między sąsiednimi węzłami redukuje się do~nierówności
\(a_x+a_y+a_z\le\tfrac12\): sumaryczna waga dyfuzyjna nie może przekraczać jednej
drugiej, co wyznacza maksymalny dopuszczalny krok czasowy:
\begin{equation}
  \Delta t_{\max}
  = \frac{1}{2\alpha\!\left(\dfrac{1}{\Delta x^2}
    + \dfrac{1}{\Delta y^2} + \dfrac{1}{\Delta z^2}\right)}.
  \label{eq:cfl}
\end{equation}
Zależność~\eqref{eq:cfl} ma istotne konsekwencje praktyczne: dwukrotne
zagęszczenie siatki wymusza czterokrotne
zmniejszenie kroku czasowego, a~zatem czterokrotny wzrost liczby kroków
potrzebnych do~przebycia tego samego przedziału czasu fizycznego. Wraz
z~ośmiokrotnym wzrostem liczby węzłów siatki daje to bardzo
szybki wzrost całkowitego nakładu obliczeniowego, jedna z~przyczyn, dla~których
symulacje dużych siatek wymagają obliczeń rozproszonych.

W~zaimplementowanym programie krok czasowy dobierany jest automatycznie
na~podstawie warunku~\eqref{eq:cfl}, z~zastosowaniem współczynnika bezpieczeństwa
\num{0.98}:
\begin{equation}
  \Delta t = \num{0.98}\cdot\Delta t_{\max}.
  \label{eq:dt}
\end{equation}
Ten \SI{2}{\percent} margines (ang.\ \emph{safety back-off}) spycha punkt pracy
poniżej granicy stabilności (\(g=\num{-0.96}\) dla~najbardziej niekorzystnego
modu), dzięki czemu schemat pozostaje po~stabilnej stronie nawet przy
uwzględnieniu błędów zaokrągleń arytmetyki zmiennoprzecinkowej i~drobnych efektów
numerycznych, które mogłyby zepchnąć krok dokładnie na~granicy w~obszar
niestabilności. Wybór
wartości bliskiej jedności jest jednocześnie korzystny wydajnościowo:
maksymalizuje postęp w~czasie fizycznym na~jeden krok obliczeniowy, minimalizując
liczbę wymaganych iteracji, a~tym samym liczbę kosztownych wymian danych między
węzłami omawianych w~dalszej części rozdziału.

\section{Zrównoleglenie obliczeń stencilowych}
\label{sec:parallel}

\subsection{Dekompozycja dziedziny i komórki brzegowe}

Obliczenia stencilowe należą do~klasy problemów zrównoleglanych przez podział
danych (ang.\ \emph{data parallelism}). Ponieważ aktualizacja
każdego węzła według~\eqref{eq:update} jest niezależna w~obrębie jednego kroku
czasowego, naturalną strategią jest dekompozycja dziedziny (ang.\ \emph{domain
decomposition}): globalna siatka zostaje podzielona na~rozłączne podobszary,
z~których każdy przypisany jest do~osobnej jednostki obliczeniowej, w~rozważanym
przypadku do~osobnego węzła klastra. Każdy węzeł oblicza aktualizację wyłącznie
dla~przydzielonych mu komórek.

Problem pojawia się na~granicach podobszarów. Aby zaktualizować komórki leżące
tuż przy krawędzi swojego podobszaru, węzeł potrzebuje wartości sąsiadów
znajdujących się już w~podobszarze sąsiedniego węzła. Rozwiązaniem jest
otoczenie każdego lokalnego podobszaru dodatkową warstwą komórek brzegowych
(ang.\ \emph{ghost cells} lub \emph{halo}), które przechowują kopie wartości
z~krawędzi sąsiednich podobszarów. W~terminologii Chapel warstwa ta nosi nazwę
\emph{fluff}; w~zaimplementowanym programie zadeklarowano ją o~grubości jednej
komórki w~każdym kierunku, \texttt{fluff=(1,1,1)}, co odpowiada zasięgowi
siedmiopunktowego stencila. \Cref{fig:domain} przedstawia schematycznie podział
dziedziny na~podobszary wraz z~warstwami brzegowymi.

\begin{figure}[htbp]
  \centering
  \begin{tikzpicture}[scale=0.95,
      core/.style={rectangle,draw=blue!55!black,thick,minimum height=20mm,minimum width=20mm,fill=blue!16},
      halo/.style={rectangle,draw=red!70!black,dashed,thick,minimum height=20mm,minimum width=4mm,fill=red!16},
      lab/.style={font=\footnotesize},
      exch/.style={{Latex[length=2mm]}-{Latex[length=2mm]},thick,codestr}]
    % three locales side by side
    \foreach \i/\x in {0/0, 1/3.2, 2/6.4} {
      \node[core] (L\i) at (\x,0) {locale \i};
    }
    % halo strips between locales (right halo of L_i / left halo of L_{i+1})
    \node[halo] (h0r) at (1.2,0) {};
    \node[halo] (h1l) at (2.0,0) {};
    \node[halo] (h1r) at (4.4,0) {};
    \node[halo] (h2l) at (5.2,0) {};
    % exchange arrows
    \draw[exch] (1.2,1.3) -- (2.0,1.3);
    \draw[exch] (4.4,1.3) -- (5.2,1.3);
    \node[lab,codestr] at (1.6,1.7) {updateFluff()};
    \node[lab,codestr] at (4.8,1.7) {updateFluff()};
    % opisy z liniami odniesienia (rozsunięte, aby się nie nakładały)
    \node[lab] (lc0) at (0,-2.5) {podobszar};
    \draw[thin,gray] (lc0.north) -- (0,-1.05);
    \node[lab] (lc1) at (3.2,-2.6) {\shortstack{obliczenia\\ lokalne}};
    \draw[thin,gray] (lc1.north) -- (3.2,-1.05);
    \node[lab] (lf) at (1.6,-3.6) {warstwa brzegowa (\emph{fluff})};
    \draw[thin,gray] (lf.north) -- (1.2,-1.05);
    % kierunek, wzdłuż którego globalna siatka jest cięta na plastry
    \draw[-{Latex[length=3mm]},thick,gray] (-1.6,-4.5) -- (8.2,-4.5);
    \node[lab,gray] at (3.2,-4.95)
      {kierunek podziału globalnej siatki na~plastry};
  \end{tikzpicture}
  \caption[Dekompozycja dziedziny z~warstwami brzegowymi]{Dekompozycja dziedziny
  w~postaci plastrów (ang.\ \emph{slab decomposition}) na~trzy jednostki
  (\emph{locale}). Globalna siatka jest cięta wzdłuż jednej osi, którą wskazuje strzałka na~dole;
  każdy podobszar, zaznaczony na~niebiesko, otoczony jest warstwą komórek brzegowych
  \emph{fluff}, zaznaczoną na~czerwono linią przerywaną, przechowującą kopie krawędzi
  sąsiadów. Przed~każdym krokiem obliczeń wywołanie \texttt{updateFluff()}
  realizuje wymianę brzegów z~najbliższymi sąsiadami.}
  \label{fig:domain}
\end{figure}

\subsection{Wymiana brzegów jako jedyna komunikacja}

W~każdym kroku czasowym warstwy brzegowe muszą zostać odświeżone, zanim rozpocznie
się faza obliczeń, gdyż przechowywane w~nich wartości pochodzą z~poprzedniego
kroku sąsiedniego węzła. Ta wymiana danych z~najbliższymi sąsiadami
(ang.\ \emph{nearest-neighbour} lub \emph{halo exchange}) jest jedynym rodzajem
komunikacji między węzłami wymaganym przez schemat jawny. W~zaimplementowanym
programie realizuje ją pojedyncze wywołanie metody \texttt{updateFluff()} biblioteki
rozkładów Chapel, wykonywane na~początku każdego kroku, przed~pętlą aktualizującą
wnętrze dziedziny.

Prowadzi to do~wyraźnego rozdziału każdego kroku czasowego na~dwie fazy:
fazę komunikacji, czyli odświeżenie warstw brzegowych, oraz fazę obliczeń, czyli wykonanie
formuły~\eqref{eq:update} dla~wszystkich komórek wnętrza. W~programie fazy te są
oddzielnie mierzone stoperami \texttt{fluffTimer} i~\texttt{computeTimer}, co
pozwala precyzyjnie rozdzielić czas spędzony na~komunikacji od~czasu obliczeń
i~stanowi podstawę analizy skalowalności w~dalszych rozdziałach. Struktura
pojedynczego kroku jest zwięźle widoczna w~procedurze \texttt{stepOnce},
przedstawionej na~\cref{lst:steponce-idea}: najpierw synchronizacja brzegów, potem
równoległa aktualizacja wnętrza.

\begin{lstlisting}[language=Chapel,caption={Pojedynczy krok czasowy: wymiana
brzegów, a~następnie równoległa aktualizacja wnętrza dziedziny.},label={lst:steponce-idea}]
proc stepOnce(ref dst, ref src, ref fluffTimer, ref computeTimer) {
  fluffTimer.start();
  src.updateFluff();               // jedyna komunikacja: wymiana halo
  fluffTimer.stop();
  computeTimer.start();
  forall (i,j,k) in interior do    // faza obliczeń, w pełni równoległa
    dst[i,j,k] = src[i,j,k]
      + ax * (src[i-1,j,k] - 2*src[i,j,k] + src[i+1,j,k])
      + ay * (src[i,j-1,k] - 2*src[i,j,k] + src[i,j+1,k])
      + az * (src[i,j,k-1] - 2*src[i,j,k] + src[i,j,k+1]);
  computeTimer.stop();
}
\end{lstlisting}

\subsection{Intensywność arytmetyczna i ograniczenia wydajności}

Kluczową cechą obliczeń stencilowych, decydującą o~ich zachowaniu wydajnościowym,
jest bardzo niska \emph{intensywność arytmetyczna} (ang.\ \emph{arithmetic
intensity})~\cite{datta-stencil}, iloraz liczby wykonanych operacji zmiennoprzecinkowych do~liczby
bajtów przeniesionych z~pamięci. Aktualizacja jednego węzła według~\eqref{eq:update}
wymaga zaledwie kilkunastu operacji zmiennoprzecinkowych, sześciu odejmowań, trzech
mnożeń przez wagi i~kilku dodawań, a~jednocześnie odczytu siedmiu wartości
i~zapisu jednej, przy \SI{8}{\byte} na~liczbę w~podwójnej precyzji daje to
zaledwie ułamek operacji na~każdy przeniesiony bajt. Co więcej, sąsiednie węzły
współdzielą większość odczytów, lecz nawet przy idealnym wykorzystaniu pamięci
podręcznej stosunek ten pozostaje niski.

Konsekwencją jest to, że wydajność obliczeń stencilowych nie jest ograniczona
przez moc obliczeniową procesora, lecz przez przepustowość transferu danych.
W~ujęciu modelu \emph{roofline}~\cite{roofline} punkt pracy tego typu zadania leży daleko
po~lewej stronie wykresu, w~obszarze ograniczonym nachyloną prostą przepustowości
pamięci, a~nie poziomą półką szczytowej wydajności obliczeniowej. Na~pojedynczym
węźle oznacza to ograniczenie przepustowością pamięci (ang.\ \emph{memory-bandwidth
bound}): dodawanie kolejnych rdzeni obliczeniowych przynosi coraz mniejsze
korzyści, gdy magistrala pamięci ulega nasyceniu. W~obliczeniach rozproszonych,
gdzie warstwy brzegowe muszą przejść przez sieć, zadanie staje się ograniczone
komunikacją (ang.\ \emph{communication bound}): o~całkowitym czasie decyduje
przepustowość i~opóźnienie połączeń sieciowych między węzłami, a~nie szybkość
samych obliczeń.

Zależność ta ma jednak korzystny aspekt związany ze~skalą problemu. Objętość
obliczeń w~podobszarze rośnie proporcjonalnie do~jego objętości, czyli jak
\(N^3\) dla~sześcianu o~boku \(N\), podczas gdy objętość danych brzegowych
podlegających wymianie rośnie proporcjonalnie do~pola powierzchni granic,
czyli jak \(N^2\). Stosunek nakładu obliczeń do~nakładu komunikacji jest zatem
rzędu \(N^3/N^2 = N\) i~rośnie liniowo z~rozmiarem krawędzi siatki. Innymi słowy,
większe sześciany mają korzystniejszy stosunek obliczeń do~komunikacji: koszt
komunikacji zostaje zamortyzowany większą liczbą obliczeń przypadających
na~każdy przesłany bajt. Obserwacja ta ma bezpośrednie przełożenie na~strategię
efektywnego skalowania i~zostanie potwierdzona empirycznie w~rozdziale
z~wynikami.

\subsection{Miary skalowalności}

Do~ilościowej oceny efektywności zrównoleglenia stosuje się standardowy aparat
pojęciowy. \emph{Przyspieszenie} (ang.\ \emph{speedup}) przy użyciu \(p\)
jednostek obliczeniowych definiuje się jako stosunek czasu wykonania sekwencyjnego
do~czasu wykonania równoległego:
\begin{equation}
  S(p) = \frac{T(1)}{T(p)},
  \label{eq:speedup}
\end{equation}
a~\emph{efektywność równoległą} (ang.\ \emph{parallel efficiency}) jako
przyspieszenie znormalizowane liczbą jednostek:
\begin{equation}
  E(p) = \frac{S(p)}{p} = \frac{T(1)}{p\,T(p)}.
  \label{eq:efficiency}
\end{equation}
Idealne, liniowe przyspieszenie odpowiada \(S(p)=p\) oraz \(E(p)=1\); w~praktyce
narzuty komunikacji, synchronizacji i~nierównomiernego obciążenia sprawiają,
że \(E(p)<1\), przy czym efektywność zwykle maleje ze~wzrostem \(p\).

Rozróżnia się dwa reżimy badania skalowalności. \emph{Skalowanie mocne}
(ang.\ \emph{strong scaling}) polega na~ustaleniu całkowitego rozmiaru problemu
i~zwiększaniu liczby jednostek obliczeniowych; celem jest skrócenie czasu
rozwiązania. W~tym reżimie udział pracy przypadającej na~jednostkę maleje,
przez co narzuty komunikacyjne stają się względnie coraz istotniejsze, a~granicę
osiągalnego przyspieszenia opisuje prawo Amdahla. \emph{Skalowanie słabe}
(ang.\ \emph{weak scaling}) polega na~ustaleniu rozmiaru problemu przypadającego
na~jedną jednostkę i~proporcjonalnym zwiększaniu zarówno liczby jednostek, jak
i~całkowitego rozmiaru problemu; celem jest rozwiązywanie coraz większych
zadań w~stałym czasie.

Fundamentalne ograniczenie skalowania mocnego formułuje \emph{prawo Amdahla}~\cite{amdahl}.
Jeśli \(f\) oznacza sekwencyjną część obliczeń, której z~natury nie da się zrównoleglić, to maksymalne osiągalne przyspieszenie przy \(p\) jednostkach
wynosi
\begin{equation}
  S(p) \le \frac{1}{f + \dfrac{1-f}{p}},
  \qquad
  \lim_{p\to\infty} S(p) = \frac{1}{f}.
  \label{eq:amdahl}
\end{equation}
Nawet niewielki udział części sekwencyjnej narzuca zatem twardą górną granicę
przyspieszenia, niezależnie od~liczby dostępnych jednostek. W~kontekście
rozważanego programu rolę nieusuwalnego narzutu pełni głównie komunikacja, wymiana warstw brzegowych, oraz synchronizacja między krokami,
co, w~połączeniu z~ograniczeniem przepustowością pamięci i~sieci omówionym
wcześniej, wyznacza realistyczne granice skalowalności analizowane w~dalszej
części pracy.

\section{Chapel, PGAS i GASNet}
\label{sec:chapel}

\subsection{Język Chapel i model PGAS}

Implementację programu oparto na~języku Chapel, języku programowania
równoległego rozwijanym z~myślą o~produktywności programisty przy zachowaniu
wysokiej wydajności na~systemach wieloprocesorowych i~rozproszonych~\cite{chapel-overview,chapel-book}. Chapel
realizuje model \emph{PGAS} (ang.\ \emph{Partitioned Global Address Space},
partycjonowana globalna przestrzeń adresowa)~\cite{pgas,upc}. Idea PGAS polega na~tym, że pamięć
wszystkich węzłów tworzy jedną logicznie wspólną przestrzeń adresową, która
jest jednak podzielona na~partycje o~jawnie określonym powinowactwie do~węzłów.
Programista operuje na~danych za~pomocą jednolitej, globalnej perspektywy
(ang.\ \emph{global-view}), tak jakby były one lokalne, natomiast środowisko
uruchomieniowe automatycznie tłumaczy dostępy do~danych zdalnych na~operacje
komunikacji sieciowej. Model ten łączy wygodę programowania w~pamięci wspólnej
z~wydajnościową świadomością rozmieszczenia danych, charakterystyczną dla~modelu
przekazywania komunikatów.

Centralnym pojęciem w~Chapel jest \emph{locale}, abstrakcja jednostki
systemu dysponującej własną pamięcią oraz zdolnością do~wykonywania obliczeń.
W~konfiguracji wielowęzłowej stosowanej w~niniejszej pracy każdemu węzłowi
fizycznemu klastra odpowiada dokładnie jeden locale, w~odwzorowaniu
\emph{jeden-do-jednego}. Wbudowana tablica \texttt{Locales} udostępnia wszystkie
dostępne jednostki, a~ich liczbę podaje stała \texttt{numLocales}. Dostęp
do~danych i~przenoszenie obliczeń pomiędzy jednostkami wyraża się wprost,
za~pomocą konstrukcji \texttt{on}.

\subsection{Iteratory równoległe i tablice globalne}

Chapel oferuje dwa uzupełniające się rodzaje równoległości. Równoległość danych
wyraża pętla \texttt{forall}, która zleca wykonanie ciała iteracji dla~wszystkich
elementów dziedziny współbieżnie, pozostawiając środowisku uruchomieniowemu
podział pracy na~zadania i~przypisanie ich do~wątków. W~programie to właśnie
pętla \texttt{forall} nad~dziedziną wnętrza (\cref{lst:steponce-idea}) realizuje fazę
obliczeń stencila. Gdy pętla iteruje po~dziedzinie rozproszonej, jej iteracje
wykonują się na~tych węzłach, do~których przypisane są odpowiednie elementy,
co zapewnia lokalność obliczeń bez~jawnego zarządzania nią przez programistę.

Równoległość zadaniową między węzłami wyraża konstrukcja
\texttt{coforall \ldots on loc}, która tworzy osobne zadanie dla~każdego locale
i~umieszcza jego wykonanie na~odpowiednim węźle. W~programie wzorzec ten
wykorzystywany jest między innymi do~równoległego zapisu klatek wynikowych, każdy węzeł niezależnie zapisuje swój lokalny podobszar. \Cref{lst:coforall}
ilustruje ten idiom.

\begin{lstlisting}[language=Chapel,caption={Równoległość zadaniowa między węzłami:
każdy locale wykonuje operację na~swoim lokalnym podobszarze danych.},label={lst:coforall}]
coforall loc in Locales do on loc {
  const ld = field.localSubdomain();   // lokalny fragment tablicy globalnej
  var block: [ld] real = field[ld];    // kopia danych bez komunikacji zdalnej
  // ... zapis / obliczenia lokalne na fragmencie ...
}
\end{lstlisting}

Tablice w~Chapel definiowane są nad~\emph{dziedzinami} (ang.\ \emph{domain}), obiektami opisującymi zbiór indeksów. Aby tablica została rozproszona pomiędzy
węzły, jej dziedzinę wyposaża się w~\emph{odwzorowanie dziedziny}
(ang.\ \emph{domain map}, dystrybucję). Najprostszą dystrybucją jest \texttt{Block},
dzieląca zakres indeksów na~w~przybliżeniu równe, przylegające bloki i~przypisująca
je kolejnym węzłom. Na~potrzeby obliczeń stencilowych Chapel udostępnia
wyspecjalizowaną dystrybucję \texttt{Stencil} z~modułu \texttt{StencilDist}, która
poza~podziałem danych utrzymuje wokół każdego lokalnego bloku warstwę komórek
brzegowych \emph{fluff}, przechowującą buforowane kopie krawędzi sąsiadów.
Warstwę tę odświeża się jawnym wywołaniem \texttt{updateFluff()}, co realizuje
opisaną wcześniej wymianę \emph{halo}~\cite{stencildist}. W~programie dziedzinę tworzy się wywołaniem
\texttt{stencilDist.\allowbreak createDomain(\ldots,\allowbreak{} fluff=(1,1,1))},
zapewniając brzeg
o~grubości jednej komórki, dokładnie odpowiadający zasięgowi
stencila~\eqref{eq:update}.

\subsection{Warstwa zadań qthreads oraz warstwa komunikacji GASNet}

Wykonanie zadań w~obrębie pojedynczego locale spoczywa na~warstwie szeregowania
zadań. W~zastosowanej konfiguracji jest to biblioteka \emph{qthreads}
(\texttt{CHPL\_TASKS=qthreads}), lekki system wątków użytkownika efektywnie
obsługujący dużą liczbę drobnoziarnistych zadań generowanych przez pętle
\texttt{forall}~\cite{qthreads}. Liczbę wątków sprzętowych dostępnych na~danym węźle udostępnia
własność \texttt{here.maxTaskPar}; wyznacza ona stopień równoległości fazy
obliczeń w~obrębie węzła i~jest raportowana przez program przy starcie.

Za~komunikację między węzłami odpowiada z~kolei warstwa \emph{GASNet}, przenośna biblioteka uruchomieniowa dostarczająca prymitywów komunikacji
jednostronnej (ang.\ \emph{one-sided}) oraz komunikatów aktywnych, na~których
Chapel opiera realizację dostępów do~pamięci zdalnej i~operacji \texttt{on}~\cite{gasnet,gasnet-spec}.
Uaktywnia się ją ustawieniem \texttt{CHPL\_COMM=gasnet}. W~zastosowanej
konfiguracji przyjęto tryb segmentu
\texttt{CHPL\_\allowbreak GASNET\_\allowbreak SEGMENT\allowbreak =everything},
w~którym cała pamięć procesu jest udostępniana dla~komunikacji zdalnej, co
upraszcza rozmieszczanie tablic rozproszonych. Warstwową strukturę całego
stosu, od~kodu aplikacji, przez model globalnej tablicy PGAS, aż po~sprzęt
sieciowy, przedstawia \cref{fig:layers}.

\begin{figure}[htbp]
  \centering
  \begin{tikzpicture}[
      layer/.style={rectangle,draw,thick,minimum width=82mm,minimum height=9mm,
                    align=center,font=\small},
      lab/.style={font=\footnotesize\itshape,gray}]
    \node[layer,fill=codekw!16] (app) {Aplikacja: program 3D (stencil, \texttt{forall}, \texttt{updateFluff})};
    \node[layer,fill=codekw!12,below=2mm of app] (chpl) {Chapel: tablice globalne PGAS, \texttt{locale}, dystrybucja \texttt{Stencil}};
    \node[layer,fill=codestr!10,below=2mm of chpl] (rt) {Warstwa uruchomieniowa: zadania \texttt{qthreads} \;$\;\vert\;$\; komunikacja \texttt{GASNet}};
    \node[layer,fill=codestr!14,below=2mm of rt] (cond) {Konduit GASNet: \texttt{udp} (AMUDP) \; / \; \texttt{mpi} (nad~MPICH/TCP)};
    \node[layer,fill=black!8,below=2mm of cond] (net) {Sieć węzłów klastra (Ethernet / TCP-IP)};
    \draw[{Latex[length=2mm]}-{Latex[length=2mm]},thick] (app) -- (chpl);
    \draw[{Latex[length=2mm]}-{Latex[length=2mm]},thick] (chpl) -- (rt);
    \draw[{Latex[length=2mm]}-{Latex[length=2mm]},thick] (rt) -- (cond);
    \draw[{Latex[length=2mm]}-{Latex[length=2mm]},thick] (cond) -- (net);
  \end{tikzpicture}
  \caption[Warstwy stosu Chapel/PGAS/GASNet]{Warstwowa architektura wykonania.
  Kod aplikacji operuje na~globalnych tablicach PGAS języka Chapel; ich dostępy
  zdalne i~wymiana warstw brzegowych są realizowane przez warstwę uruchomieniową,
  w~której zadania \texttt{qthreads} wykonują się lokalnie, a~komunikacja \texttt{GASNet}
  zachodzi między~węzłami; warstwa ta korzysta z~kolei z~wybranego konduitu
  odwzorowującego operacje na~fizyczną sieć klastra.}
  \label{fig:layers}
\end{figure}

\subsection{Konduity GASNet}

GASNet abstrahuje sprzęt sieciowy poprzez wymienne \emph{konduity}
(ang.\ \emph{conduits}), moduły implementujące prymitywy komunikacji
na~konkretnej technologii transportu. W~ramach niniejszej pracy porównywano dwa
konduity dostępne na~klastrze zbudowanym z~węzłów połączonych zwykłą siecią
Ethernet. Konduit \texttt{udp}, oparty na~warstwie AMUDP, realizuje komunikaty
aktywne bezpośrednio nad~datagramami UDP i~jest uruchamiany za~pomocą programu
rozruchowego \texttt{amudprun}. Jego zaletą jest prostota konfiguracji, wadą zaś
podatność na~przeciążenie przy intensywnej, drobnoziarnistej wymianie danych.
Konduit \texttt{mpi} przenosi komunikację GASNet na~istniejącą implementację
MPI, w~zastosowanej konfiguracji opartej na~MPICH nad~TCP, i~jest uruchamiany
programem \texttt{gasnetrun\_mpi}; zapewnia on stabilniejszy i~wydajniejszy
transport dla~podtrzymywanych, regularnych wzorców komunikacji, jakie generuje
cykliczna wymiana warstw brzegowych w~opisywanym programie.

Wybór konduitu okazuje się mieć istotny wpływ na~zachowanie rozproszonej
symulacji, ponieważ, jak wykazano w~\cref{sec:parallel}, w~reżimie
wielowęzłowym zadanie jest ograniczone komunikacją, a~jedyną komunikacją
per~krok jest wymiana \emph{halo} realizowana przez \texttt{updateFluff()}.
To właśnie właściwości warstwy komunikacyjnej, jej przepustowość, opóźnienie
oraz odporność na~przeciążenie, decydują w~praktyce o~osiąganej skalowalności,
co czyni dobór konduitu jednym z~centralnych zagadnień badawczych analizowanych
w~dalszej części pracy.
